\documentclass[a4paper,11pt]{article}

\usepackage[utf8]{inputenc}
\usepackage[T1]{fontenc}
\usepackage[spanish,es-nodecimaldot,es-noshorthands]{babel}
\usepackage[top=2.2cm,bottom=2.2cm,left=2.4cm,right=2.4cm,headheight=14pt]{geometry}
\usepackage{amsmath,amssymb,amsthm}
\usepackage{booktabs,array}
\usepackage{enumitem}
\usepackage{xcolor}
\usepackage[most]{tcolorbox}
\usepackage{fancyhdr}
\usepackage{listings}
\usepackage{hyperref}

\definecolor{azul}{HTML}{174A70}
\definecolor{fondo}{HTML}{F3F6FA}
\hypersetup{colorlinks=true,linkcolor=azul,urlcolor=azul}
\setlength{\parindent}{0pt}
\setlength{\parskip}{5pt}
\setlist{itemsep=3pt}
\pagestyle{fancy}
\fancyhf{}
\lhead{\textsc{Investigación de Operaciones}}
\chead{\small Unidad II}
\rhead{\thepage}
\renewcommand{\headrulewidth}{0.4pt}

\newtcolorbox{idea}[1][Idea clave]{enhanced,breakable,colback=fondo,colframe=azul,boxrule=0.5pt,arc=2pt,
  fonttitle=\bfseries\small,title={#1}}
\theoremstyle{plain}
\newtheorem{teorema}{Teorema}[section]
\theoremstyle{definition}
\newtheorem{definicion}[teorema]{Definición}
\newtheorem{ejemplo}[teorema]{Ejemplo}
\renewcommand{\qedsymbol}{$\blacksquare$}

\lstset{basicstyle=\ttfamily\small,backgroundcolor=\color{fondo},frame=single,rulecolor=\color{azul!40},
  columns=fullflexible,keepspaces=true,literate={á}{{\'a}}1{é}{{\'e}}1{í}{{\'i}}1{ó}{{\'o}}1{ú}{{\'u}}1{ñ}{{\~n}}1}

\newcommand{\fr}[2]{\tfrac{#1}{#2}}

\title{\textcolor{azul}{\Huge Investigación de Operaciones}\\[6pt]
\Large Unidad II: dualidad, precios sombra y análisis de sensibilidad\\[4pt]
\large Apuntes guiados, ejemplos y ejercicios}
\author{Julio César Galindo López}
\date{}

\begin{document}
\maketitle

\section*{Motivación}
En la Unidad I aprendimos a responder la pregunta \emph{¿cuánto producir?} La solución óptima, sin
embargo, contiene una segunda respuesta que la empresa necesita tanto como la primera: \emph{¿cuánto vale
cada recurso?} Si una hora adicional de máquina permite ganar más, ¿cuánto conviene pagar por ella? Si un
proveedor ofrece materia prima extra, ¿a qué precio deja de ser buen negocio? Si se propone un producto
nuevo, ¿vale la pena fabricarlo sin resolver todo el problema otra vez?

Las tres preguntas tienen la misma respuesta matemática: cada programa lineal tiene asociado otro programa
lineal, su \emph{dual}, cuyas variables son precios de los recursos. Esta unidad construye el dual, demuestra
por qué su valor óptimo coincide con el del problema original y enseña a leer en la tabla óptima del
simplex los precios sombra y los rangos dentro de los cuales la solución sigue siendo válida.

\section*{Cómo estudiar esta unidad}
Toda la unidad gira alrededor del ejemplo de producción resuelto en la Unidad I, de modo que cada concepto
nuevo se aplica a números ya conocidos. Se recomienda tener a la vista su tabla óptima. Las secciones 2 y 3
construyen el dual; la 4 demuestra los teoremas que lo justifican; las 5 y 6 interpretan la solución; la 7
estudia la sensibilidad, y la 8 muestra cómo comprobar todo en Python.

\tableofcontents
\newpage

%=====================================================================
\section{El problema de referencia}

Una empresa fabrica dos productos, A y B, con utilidades de $4$ y $6$ (miles de pesos por lote). Cada lote
de A requiere $0.5$ horas de máquina y $12.5$ kg de materia prima; cada lote de B, $1$ hora y $10$ kg. Se
dispone de $600$ horas y $10\,000$ kg, y el mercado absorbe a lo más $700$ lotes de A. El modelo es
\[
\begin{aligned}
\max\ Z&=4x_1+6x_2\\
\text{s.a.}\quad 0.5x_1+x_2&\le600 &&\text{(horas de máquina)}\\
12.5x_1+10x_2&\le10\,000 &&\text{(materia prima)}\\
x_1&\le700 &&\text{(mercado de A)}\\
x_1,x_2&\ge0.
\end{aligned}
\]
En la Unidad I el simplex terminó con las variables básicas $x_1,x_2,s_3$ y
\[
x_1=\fr{1600}{3}\approx533.33,\qquad x_2=\fr{1000}{3}\approx333.33,\qquad s_3=\fr{500}{3}\approx166.67,
\qquad s_1=s_2=0,\qquad Z^*=\fr{12\,400}{3}\approx4133.33.
\]
Las horas y la materia prima se agotan; sobra capacidad de mercado de A.

%=====================================================================
\section{El problema dual}

\subsection{La pregunta del comprador}
Imagine que un competidor quiere comprar \emph{todos} los recursos de la empresa: las $600$ horas, los
$10\,000$ kg y los $700$ ``derechos de mercado'' de A. Ofrece un precio $y_1$ por hora, $y_2$ por kilogramo y
$y_3$ por derecho. ¿Qué precios mínimos debe ofrecer para que la empresa acepte vender en lugar de
producir?

Un lote de A usa $0.5$ horas, $12.5$ kg y un derecho de mercado, y deja $4$ de utilidad. La empresa solo
vende si lo que recibe por esos recursos es al menos lo que ganaría produciendo:
$0.5y_1+12.5y_2+y_3\ge4$. Lo mismo con B: $y_1+10y_2\ge6$. El comprador quiere pagar lo menos posible:
\[
\begin{aligned}
\min\ W&=600y_1+10\,000y_2+700y_3\\
\text{s.a.}\quad 0.5y_1+12.5y_2+y_3&\ge4 &&\text{(producto A)}\\
y_1+10y_2&\ge6 &&\text{(producto B)}\\
y_1,y_2,y_3&\ge0.
\end{aligned}
\]
Este es el \textbf{dual} del problema de producción, que llamamos \textbf{primal}.

\subsection{Reglas de construcción}

\begin{definicion}[Par primal--dual en forma canónica]
Si el primal es $\max\ \mathbf c^\top\mathbf x$ sujeto a $A\mathbf x\le\mathbf b$, $\mathbf x\ge\mathbf 0$, su dual es
$\min\ \mathbf b^\top\mathbf y$ sujeto a $A^\top\mathbf y\ge\mathbf c$, $\mathbf y\ge\mathbf 0$.
\end{definicion}

En palabras: la matriz se transpone, los lados derechos del primal pasan a ser los coeficientes del objetivo
dual y viceversa, y hay una variable dual por cada restricción del primal.

\begin{center}
\begin{tabular}{@{}ll@{}}\toprule
Primal (máx.) & Dual (mín.)\\\midrule
restricción $i$ de tipo $\le$ & variable $y_i\ge0$\\
restricción $i$ de tipo $\ge$ & variable $y_i\le0$\\
restricción $i$ de tipo $=$ & variable $y_i$ libre\\
variable $x_j\ge0$ & restricción $j$ de tipo $\ge$\\
variable $x_j\le0$ & restricción $j$ de tipo $\le$\\
variable $x_j$ libre & restricción $j$ de tipo $=$\\
\bottomrule
\end{tabular}
\end{center}

\begin{idea}
El dual del dual es el primal. Por eso da lo mismo cuál de los dos llamamos primal: la tabla se lee en
ambos sentidos.
\end{idea}

\begin{ejemplo}[Dual de un problema de minimización]
El problema de dieta $\min\ 3x_1+2x_2$ s.a.\ $x_1+x_2\ge4$, $2x_1+x_2\ge6$, $x_1,x_2\ge0$ tiene por dual
$\max\ 4y_1+6y_2$ s.a.\ $y_1+2y_2\le3$, $y_1+y_2\le2$, $y_1,y_2\ge0$. Aquí $y_i$ es el precio que un
fabricante de suplementos podría cobrar por cada unidad del nutriente $i$ sin que convenga comprar los
alimentos.
\end{ejemplo}

\begin{ejemplo}[Restricción de igualdad]
Si al problema de referencia se le exige usar exactamente las $600$ horas, la primera restricción es
$0.5x_1+x_2=600$ y la variable dual $y_1$ queda libre de signo: una hora puede tener valor negativo si
obliga a producir algo que no conviene.
\end{ejemplo}

%=====================================================================
\section{Resolver el dual del ejemplo}

El dual tiene tres variables, pero la Sección 4 mostrará que podemos adivinar su solución a partir del
primal. Como sobra mercado de A ($s_3>0$), ese recurso no tiene valor marginal: $y_3=0$. Como $x_1>0$ y
$x_2>0$, las dos restricciones duales se cumplen con igualdad:
\[
0.5y_1+12.5y_2=4,\qquad y_1+10y_2=6.
\]
De la segunda, $y_1=6-10y_2$; sustituyendo, $3-5y_2+12.5y_2=4$, es decir $7.5y_2=1$. Entonces
\[
y_2=\fr{2}{15}\approx0.133,\qquad y_1=6-\fr{4}{3}=\fr{14}{3}\approx4.667,\qquad y_3=0,
\]
y
\[
W^*=600\cdot\fr{14}{3}+10\,000\cdot\fr{2}{15}+0=2800+\fr{4000}{3}=\fr{12\,400}{3}=Z^*.
\]
Los dos óptimos coinciden. No es casualidad.

%=====================================================================
\section{Los teoremas de dualidad}

En lo que sigue el primal es $\max\ \mathbf c^\top\mathbf x$ s.a.\ $A\mathbf x\le\mathbf b$, $\mathbf x\ge0$, y el dual es
$\min\ \mathbf b^\top\mathbf y$ s.a.\ $A^\top\mathbf y\ge\mathbf c$, $\mathbf y\ge0$.

\begin{teorema}[Dualidad débil]
Si $\mathbf x$ es factible para el primal y $\mathbf y$ es factible para el dual, entonces
$\mathbf c^\top\mathbf x\le\mathbf b^\top\mathbf y$.
\end{teorema}
\begin{proof}
Como $\mathbf x\ge0$ y $A^\top\mathbf y\ge\mathbf c$, se tiene $\mathbf c^\top\mathbf x\le(A^\top\mathbf y)^\top\mathbf x=\mathbf y^\top A\mathbf x$. Como
$\mathbf y\ge0$ y $A\mathbf x\le\mathbf b$, se tiene $\mathbf y^\top A\mathbf x\le\mathbf y^\top\mathbf b$.
\end{proof}

Consecuencias inmediatas:
\begin{enumerate}
\item Cualquier solución factible del dual da una \emph{cota superior} para la utilidad máxima, sin resolver
el primal. En el ejemplo, $\mathbf y=(6,0,4)$ es factible y da $W=600\cdot6+700\cdot4=6400$: la utilidad
nunca puede pasar de $6400$.
\item Si $\mathbf c^\top\mathbf x=\mathbf b^\top\mathbf y$ para un par factible, ambos son óptimos: ninguna $\mathbf x$ puede dar más
que esa cota, y ninguna $\mathbf y$ puede dar menos.
\item Si el primal es no acotado, el dual es infactible (no existe ninguna cota superior), y viceversa.
\end{enumerate}

\begin{teorema}[Dualidad fuerte]
Si el primal tiene solución óptima $\mathbf x^*$, el dual también tiene solución óptima $\mathbf y^*$ y
$\mathbf c^\top\mathbf x^*=\mathbf b^\top\mathbf y^*$.
\end{teorema}
\begin{proof}[Idea de la demostración]
Sea $B$ la matriz básica de la tabla óptima del simplex y $\mathbf c_B$ los costos de las variables básicas.
Defina $\mathbf y^{*\top}=\mathbf c_B^\top B^{-1}$. La optimalidad ($C_j-Z_j\le0$ para todas las columnas) dice
exactamente que $A^\top\mathbf y^*\ge\mathbf c$ (columnas de las variables de decisión) y que $\mathbf y^*\ge0$
(columnas de las holguras). Así $\mathbf y^*$ es factible para el dual y
$\mathbf b^\top\mathbf y^*=\mathbf c_B^\top B^{-1}\mathbf b=\mathbf c_B^\top\mathbf x_B=Z^*$; por la dualidad débil es óptima.
\end{proof}

\begin{teorema}[Holgura complementaria]
Sean $\mathbf x$ y $\mathbf y$ factibles. Son óptimas si y solo si
\[
y_i\,(b_i-(A\mathbf x)_i)=0\ \ \text{para toda } i,\qquad x_j\,((A^\top\mathbf y)_j-c_j)=0\ \ \text{para toda } j.
\]
\end{teorema}
\begin{proof}
En la demostración de la dualidad débil,
$\mathbf b^\top\mathbf y-\mathbf c^\top\mathbf x=\mathbf y^\top(\mathbf b-A\mathbf x)+\mathbf x^\top(A^\top\mathbf y-\mathbf c)$, suma de términos no negativos.
Es cero, es decir hay optimalidad, si y solo si cada término es cero.
\end{proof}

\begin{idea}[Lectura económica de la holgura complementaria]
\begin{itemize}
\item Un recurso que sobra tiene precio sombra cero: nadie paga por más de algo que ya sobra ($s_3>0\Rightarrow y_3=0$).
\item Un recurso con precio positivo se agota ($y_1>0\Rightarrow s_1=0$).
\item Un producto que se fabrica cubre exactamente el valor de los recursos que consume
($x_1>0\Rightarrow0.5y_1+12.5y_2+y_3=4$).
\item Un producto cuyo costo de oportunidad supera su utilidad no se fabrica.
\end{itemize}
\end{idea}

%=====================================================================
\section{El dual está escrito en la tabla óptima}

No hace falta resolver el dual por separado: el simplex lo resuelve al mismo tiempo. Con la convención
$C_j-Z_j$ de la Unidad I, en la tabla óptima de maximización
\[
y_i=Z_j\ \text{de la columna de la holgura } s_i,\qquad\text{es decir,}\qquad y_i=-(C_j-Z_j)_{s_i}.
\]

Para el ejemplo, la matriz básica (columnas de $x_1$, $x_2$ y $s_3$ en las tres restricciones) y su
inversa son
\[
B=\begin{pmatrix}0.5&1&0\\12.5&10&0\\1&0&1\end{pmatrix},\qquad
B^{-1}=\begin{pmatrix}-\fr43&\fr2{15}&0\\[2pt]\fr53&-\fr1{15}&0\\[2pt]\fr43&-\fr2{15}&1\end{pmatrix}.
\]
Las columnas de $s_1$ y $s_2$ en la tabla óptima son las dos primeras columnas de $B^{-1}$, y
\[
\mathbf y^{*\top}=\mathbf c_B^\top B^{-1}=(4,\ 6,\ 0)\,B^{-1}=\Big(\fr{14}{3},\ \fr{2}{15},\ 0\Big),
\qquad
\mathbf x_B=B^{-1}\mathbf b=\Big(\fr{1600}{3},\ \fr{1000}{3},\ \fr{500}{3}\Big).
\]

\begin{center}
\begin{tabular}{@{}c|rrrrr|r@{}}\toprule
Básica & $x_1$ & $x_2$ & $s_1$ & $s_2$ & $s_3$ & $b$\\\midrule
$x_1$ & 1 & 0 & $-\fr43$ & $\fr2{15}$ & 0 & $\fr{1600}3$\\[2pt]
$x_2$ & 0 & 1 & $\fr53$ & $-\fr1{15}$ & 0 & $\fr{1000}3$\\[2pt]
$s_3$ & 0 & 0 & $\fr43$ & $-\fr2{15}$ & 1 & $\fr{500}3$\\\midrule
$Z_j$ & 4 & 6 & $\fr{14}3$ & $\fr2{15}$ & 0 & $\fr{12400}3$\\[2pt]
$C_j-Z_j$ & 0 & 0 & $-\fr{14}3$ & $-\fr2{15}$ & 0 & \\\bottomrule
\end{tabular}
\end{center}

La fila $Z_j$ bajo las holguras contiene los precios sombra $\big(\fr{14}{3},\fr{2}{15},0\big)$, que coinciden con
la solución del dual obtenida en la Sección 3.

%=====================================================================
\section{Precios sombra}

\begin{definicion}
El \emph{precio sombra} del recurso $i$ es el aumento del valor óptimo del objetivo por cada unidad
adicional del recurso $i$, mientras la base óptima no cambie. Es el valor óptimo de la variable dual $y_i$.
\end{definicion}

En el ejemplo:
\begin{itemize}
\item \textbf{Horas de máquina:} $y_1=\fr{14}3\approx4.67$ miles de pesos por hora. La empresa debería estar
dispuesta a pagar hasta $4\,667$ pesos por una hora extra (por ejemplo, un turno adicional).
\item \textbf{Materia prima:} $y_2=\fr2{15}\approx0.133$ miles de pesos por kg, es decir $133$ pesos por kg
adicional. Si un proveedor ofrece materia prima extra a menos de eso, conviene comprarla.
\item \textbf{Mercado de A:} $y_3=0$. Ampliar el mercado de A no sirve hoy: ya sobran $166.67$ lotes.
\end{itemize}

\begin{ejemplo}[Verificación directa]
Con $601$ horas, las restricciones activas $0.5x_1+x_2=601$ y $12.5x_1+10x_2=10\,000$ dan
$x_1=532$ y $x_2=335$, con utilidad $4\cdot532+6\cdot335=4138$; el aumento es
$4138-\fr{12\,400}3=\fr{14}3$. La primera columna de $B^{-1}$ explica el mecanismo: una hora más cambia
$(x_1,x_2,s_3)$ en $\big(-\fr43,\fr53,\fr43\big)$. La hora extra se usa para producir más B y menos A,
liberando materia prima.
\end{ejemplo}

\begin{idea}[Cuidado]
El precio sombra es un valor \emph{marginal} y local. No dice cuánto vale comprar $1000$ horas más: solo es
válido mientras la base óptima no cambie. La Sección 7 calcula hasta dónde.
\end{idea}

%=====================================================================
\section{Costo reducido y productos nuevos}

\begin{definicion}
El \emph{costo de oportunidad} de una actividad $j$ es $\sum_i a_{ij}y_i$: el valor, a precios sombra, de los
recursos que consume. Su \emph{costo reducido} es $c_j-\sum_i a_{ij}y_i$ (en maximización).
\end{definicion}

Una actividad que no está en la solución solo conviene si su costo reducido es positivo. Esto permite
evaluar un producto nuevo \emph{sin} volver a resolver el problema.

\begin{ejemplo}[¿Lanzar el producto C?]
Se propone un producto C que requiere $0.8$ horas y $15$ kg por lote y deja $7$ de utilidad. Su costo de
oportunidad es
\[
0.8\cdot\fr{14}{3}+15\cdot\fr{2}{15}=3.733+2=5.733.
\]
Como $7>5.733$, el costo reducido es $7-5.733=1.267>0$: conviene producir C. Al resolver de nuevo con C el
óptimo es $x_1=0$, $x_2=\fr{1000}{7}\approx142.86$, $x_3=\fr{4000}{7}\approx571.43$, con
$Z^*=\fr{34\,000}{7}\approx4857.14$. Si C dejara $5$ de utilidad, la prueba diría que no conviene, y la
solución actual seguiría siendo óptima.
\end{ejemplo}

%=====================================================================
\section{Análisis de sensibilidad}

Los datos de un modelo casi nunca son exactos: utilidades, consumos y disponibilidades se estiman. El
análisis de sensibilidad pregunta cuánto puede cambiar un dato sin que cambie la base óptima.

\subsection{Coeficientes de la función objetivo}
La solución óptima es el vértice donde se cruzan las rectas de horas ($x_2=600-0.5x_1$, pendiente $-0.5$) y
de materia prima ($x_2=1000-1.25x_1$, pendiente $-1.25$). Ese vértice sigue siendo óptimo mientras la
pendiente de las rectas de isoutilidad, $-c_1/c_2$, quede entre las dos:
\[
0.5\le\frac{c_1}{c_2}\le1.25.
\]
\begin{itemize}
\item Con $c_2=6$ fijo: $3\le c_1\le7.5$. Si la utilidad de A baja de $3$, conviene producir solo B
($x_1=0$, $x_2=600$); si sube de $7.5$, conviene llegar al tope de mercado de A ($x_1=700$, $x_2=125$).
\item Con $c_1=4$ fijo: $3.2\le c_2\le8$.
\end{itemize}
Dentro de estos rangos la \emph{mezcla} de producción no cambia, pero la utilidad total sí.

\subsection{Lados derechos}
Si cambia $b_i$, la nueva solución básica es $\mathbf x_B=B^{-1}\mathbf b$; la base sigue siendo óptima mientras
$\mathbf x_B\ge0$ (los precios sombra no dependen de $\mathbf b$). Para las horas, con $b_1=600+\Delta$:
\[
\mathbf x_B=\Big(\fr{1600}3-\fr43\Delta,\ \ \fr{1000}3+\fr53\Delta,\ \ \fr{500}3+\fr43\Delta\Big)\ge0
\iff -125\le\Delta\le400.
\]
Es decir, el precio sombra $\fr{14}3$ es válido para $475\le b_1\le1000$ horas. Del mismo modo:
\begin{center}
\begin{tabular}{@{}lccc@{}}\toprule
Recurso & Disponible & Precio sombra & Rango de validez\\\midrule
Horas de máquina & $600$ & $\fr{14}3\approx4.667$ & $[475,\ 1000]$\\
Materia prima (kg) & $10\,000$ & $\fr2{15}\approx0.133$ & $[6000,\ 11\,250]$\\
Mercado de A & $700$ & $0$ & $[533.33,\ \infty)$\\
\bottomrule
\end{tabular}
\end{center}

\begin{ejemplo}
Un turno extra de $200$ horas cuesta $600$ (miles de pesos). Como $800$ está en el rango, la utilidad
aumenta $200\cdot\fr{14}3=933.33$, más que el costo: el turno conviene, con ganancia neta $333.33$. Un turno de
$500$ horas llevaría a $1100>1000$ y el precio sombra ya no aplica: hay que resolver de nuevo.
\end{ejemplo}

\subsection{Una restricción nueva}
Si se agrega una restricción, primero se verifica si la solución actual la cumple. Si la cumple, sigue
siendo óptima (agregar restricciones no puede mejorar el óptimo). Si no la cumple, la solución cambia y se
reoptimiza (por ejemplo, con el simplex dual).

\begin{ejemplo}
Un límite de almacén $x_1+x_2\le900$: la solución actual da $533.33+333.33=866.67\le900$; sigue siendo
óptima. Con un límite de $800$, la solución actual lo viola y hay que reoptimizar.
\end{ejemplo}

%=====================================================================
\section{Comprobación en Python}

\texttt{scipy.optimize.linprog} resuelve el primal y reporta los precios sombra en
\texttt{ineqlin.marginals} (con signo negativo, porque \texttt{linprog} minimiza):

\begin{lstlisting}[language=Python]
from scipy.optimize import linprog

c = [-4, -6]                       # linprog minimiza: se cambia el signo
A = [[0.5, 1], [12.5, 10], [1, 0]]
b = [600, 10000, 700]
res = linprog(c, A_ub=A, b_ub=b, method="highs")
print(res.x, -res.fun)             # [533.33 333.33] 4133.33
print(-res.ineqlin.marginals)      # precios sombra: [4.667 0.133 0.   ]
print(res.ineqlin.residual)        # holguras: [0. 0. 166.67]

# el dual: min b'y  s.a. A'y >= c, y >= 0
import numpy as np
dual = linprog(b, A_ub=-np.array(A).T, b_ub=[-4, -6], method="highs")
print(dual.x, dual.fun)            # [4.667 0.133 0.] 4133.33
\end{lstlisting}

Para comprobar un rango basta repetir el cálculo con el dato modificado: con $b_1=1001$ el precio sombra
de las horas cae a cero, porque A deja de producirse y la base cambia.

%=====================================================================
\section{Errores frecuentes}
\begin{itemize}
\item Leer el precio sombra como el \emph{precio de mercado} del recurso. Es su valor para esta empresa, con
este plan, en el margen.
\item Usar el precio sombra fuera de su rango de validez.
\item Olvidar que en minimización los signos de las variables duales se invierten según el tipo de restricción.
\item Concluir que un recurso con precio sombra cero no sirve: sirve, pero hoy sobra.
\item Redondear la solución continua y suponer que sigue siendo óptima o factible.
\end{itemize}

%=====================================================================
\section{Ejercicios resueltos}

\textbf{1.} Escriba el dual de $\max\ 3x_1+5x_2$ s.a.\ $x_1\le4$, $2x_2\le12$, $3x_1+2x_2\le18$, $x_1,x_2\ge0$.

\emph{Solución.} $\min\ 4y_1+12y_2+18y_3$ s.a.\ $y_1+3y_3\ge3$, $2y_2+2y_3\ge5$, $y\ge0$. El primal tiene
óptimo $x=(2,6)$, $Z=36$; por holgura complementaria, como sobra la primera restricción ($2<4$), $y_1=0$, y
de $3y_3=3$, $2y_2+2y_3=5$ se obtiene $y=(0,\fr32,1)$ con $W=12\cdot\fr32+18=36$.

\textbf{2.} En el problema de referencia, ¿cuánto pagaría como máximo por $500$ kg extra de materia prima?

\emph{Solución.} $10\,500$ kg está en $[6000,11\,250]$, así que la utilidad sube $500\cdot\fr2{15}=66.67$ miles de
pesos. Es el precio máximo por el lote completo, $133.33$ pesos por kg.

\textbf{3.} ¿Sigue siendo óptima la mezcla actual si la utilidad de B sube a $7.5$?

\emph{Solución.} Sí: $7.5\in[3.2,8]$. La nueva utilidad es $4\cdot\fr{1600}3+7.5\cdot\fr{1000}3=\fr{13\,900}3\approx4633.33$.

\textbf{4.} Verifique la holgura complementaria en la solución del ejemplo.

\emph{Solución.} $y_1s_1=\fr{14}3\cdot0=0$, $y_2s_2=0$, $y_3s_3=0\cdot\fr{500}3=0$; las dos restricciones duales
se cumplen con igualdad y $x_1,x_2>0$. Todos los productos son cero.

%=====================================================================
\section{Ejercicios propuestos}
\begin{enumerate}
\item Escriba el dual del problema de minimización con artificiales de la Unidad I y resuélvalo por holgura complementaria.
\item En el problema de referencia, encuentre el rango de validez del precio sombra de la materia prima y explique qué ocurre en cada extremo.
\item Un producto D usa $0.4$ horas y $14$ kg y deja $3.5$. ¿Conviene fabricarlo? Justifique con su costo reducido y compruebe en Python.
\item Muestre con un ejemplo de dos variables que si el primal es no acotado el dual es infactible.
\item Demuestre que el dual del dual es el primal.
\item El dueño propone un límite de calidad $x_2\le300$. ¿Cambia la solución? Si sí, reoptimice e interprete el nuevo precio sombra de las horas.
\item Construya un problema en el que un recurso agotado tenga precio sombra cero (degeneración) y explique por qué no contradice la holgura complementaria.
\end{enumerate}

\section*{Lista de verificación al terminar la unidad}
\begin{enumerate}
\item ¿Escribo el dual de cualquier programa lineal usando la tabla de reglas?
\item ¿Demuestro la dualidad débil y explico por qué da cotas sin resolver nada?
\item ¿Leo los precios sombra en la tabla óptima y los interpreto con unidades?
\item ¿Uso la holgura complementaria para obtener una solución a partir de la otra?
\item ¿Evalúo un producto nuevo con su costo reducido?
\item ¿Calculo rangos de validez para coeficientes del objetivo y lados derechos?
\item ¿Compruebo todo con \texttt{linprog}?
\end{enumerate}

\section*{Referencias}
\begin{itemize}
\item Hillier, F.\,S. y Lieberman, G.\,J. (2021). \emph{Introduction to Operations Research}, 11.ª ed. McGraw-Hill. Caps.~6 y~7.
\item Taha, H.\,A. (2017). \emph{Operations Research: An Introduction}, 10.ª ed. Pearson. Cap.~4.
\item Bertsimas, D. y Tsitsiklis, J.\,N. (1997). \emph{Introduction to Linear Optimization}. Athena Scientific. Cap.~4.
\end{itemize}

\end{document}
